\solchapter{strings}

\sol{strings}{1} Write $\hbar^xc^yG^z$, with $[\hbar] = \mathrm{ML^2T^{-1}}$, $[c] = \mathrm{LT^{-1}}$, $[G] = \mathrm{M^{-1}L^3T^{-2}}$. For a length, the powers of M, L, and T must be $0$, $1$, and $0$: $x - z = 0$, $2x+y+3z = 1$, $-x-y-2z = 0$, with the unique solution $x = z = \tfrac12$, $y = -\tfrac32$. Similarly for mass and time. Hence $\ell_P = \sqrt{\hbar G/c^3} = 1.6\times10^{-35}\un{m}$, $m_P = \sqrt{\hbar c/G} = 2.2\times10^{-8}\un{kg}$ (about a flea's egg), $t_P = \sqrt{\hbar G/c^5} = 5.4\times10^{-44}\un{s}$, and $T_P = m_Pc^2/\kB = 1.4\times10^{32}\un{K}$.

\sol{strings}{2} $E_P = 1.2\times10^{19}\un{GeV} = 1.2\times10^{28}\un{eV}$: $15$ orders of magnitude above the LHC, and $8$ above the most energetic cosmic rays.

\sol{strings}{3} $\hbar/mc = 2Gm/c^2$ gives $m = \sqrt{\hbar c/2G} = m_P/\sqrt2 = 1.5\times10^{-8}\un{kg}$. To localize a particle within $\Delta x$, you need a probe of wavelength $\Delta x$, carrying energy $\sim\hbar c/\Delta x$. Below the Planck length, that energy in that region forms a black hole larger than $\Delta x$: the harder you try to look, the bigger the black hole that hides what you are looking at. Distances below $\ell_P$ may not mean anything.

\sol{strings}{4} $S = \kB A/4\ell_P^2 = \pi\kB r_s^2/\ell_P^2$, with $r_s = 3.0\un{km}$: $S = 1.0\times10^{77}\kB$, or $1.5\times10^{77}$ bits --- $10^{19}$ times the entropy of the Sun. Since $S\propto M^2$, Sagittarius~A$^*$ ($4\times10^6\Msun$) has $S\approx2\times10^{90}\kB$, more than the entire cosmic microwave background. The entropy of the universe is dominated by its black holes.

\sol{strings}{5} $r_s = 2GM/c^2 = 3\times10^{-16}\un{m}$, about a third of a proton's radius. $T = \hbar c^3/8\pi GM\kB = 6\times10^{11}\un{K}$, or $\kB T\approx50\un{MeV}$: it emits not only photons but electrons, muons, and pions. As it shrinks it heats up and radiates faster, and the end is an explosion. Since lifetime scales as $M^3$, the black hole in its last second has $2.6\times10^5\un{kg}$, releasing $2\times10^{22}\un{J}$ --- about five million megatons --- in a final burst of gamma rays. Telescopes have looked for such bursts and seen none, which limits how many primordial black holes there could be.

\sol{strings}{6} A massless particle in five dimensions has $E^2 = (p^2 + p_5^2)c^2$. Seen in four dimensions, $p_5c$ plays the role of a rest energy $mc^2$, so $m_n = \abs{p_5}/c = n\hbar/Rc$: a tower of equally spaced masses. No tower below a few TeV means $\hbar c/R\gtrsim3\un{TeV}$, so $R\lesssim\hbar c/(3\un{TeV}) = 197\un{MeV\,fm}/(3\times10^6\un{MeV})\approx7\times10^{-20}\un{m}$: a ten-thousandth of a proton.

\sol{strings}{7} The gravitational energy is $G(E/c^2)^2/r$ with $r = \hbar c/E$, i.e.\ $GE^3/\hbar c^5$; dividing by $E$ gives $GE^2/\hbar c^5 = (E/E_P)^2$, using $E_P^2 = \hbar c^5/G$. Two electrons at rest: $(0.511\un{MeV}/E_P)^2 = 2\times10^{-45}$. Two LHC protons: $(6.8\un{TeV}/E_P)^2 = 3\times10^{-31}$. Gravity is weak, but like Fermi's theory, its effective strength grows as $E^2$ and reaches $1$ at $E_P$, where perturbation theory breaks down. Fermi's theory broke down at $300\un{GeV}$ and the $W$ appeared; something new must appear by the Planck energy --- perhaps strings.

\sol{strings}{8} $\pi r^2/\ell_P^2 = 1.2\times10^{70}$ for $r = 1\un{m}$. A sphere of air of that radius ($4\un{m^3}$) has about $4\times10^{26}\kB$: $10^{43}$ below the bound. Ordinary matter has entropy proportional to its volume, roughly one $\kB$ per particle. To approach the bound you would have to pack in so much matter that it would collapse into a black hole first. Only black holes saturate the bound.

\sol{strings}{9} $\rho = p/0.3B = 1.2\times10^{19}/(0.3\times8)\un{m} = 5\times10^{18}\un{m}$, about $540$ light-years, comparable to the distances of many of the stars you can see with the naked eye. The synchrotron loss per turn, $88.5\,E^4/\rho\times(m_e/m_p)^4$ in keV, comes to $3\times10^{40}\un{GeV}$ --- $10^{21}$ times the beam's energy. The protons would radiate their energy away almost at once. Even in principle, there is no escape: concentrating the Planck energy in a Planck-sized region makes a Planck-mass black hole. Colliding particles at still higher energies makes bigger black holes, which probe \emph{larger} distances, not smaller. Gravity itself seems to forbid seeing below the Planck length.

\sol{strings}{10} $\hbar/4\pi\kB = 6.1\times10^{-13}\un{K\,s}$. Water: $10^{-3}/(4\times10^6) = 2.5\times10^{-10}\un{K\,s}$, about $400$ times the bound. In kinetic theory, $\eta\sim np\ell$ (momentum carried across a distance $\ell$, the mean free path) and $s\sim n\kB$, so $\eta/s\sim p\ell/\kB$; the uncertainty principle requires $p\ell\gtrsim\hbar$. A fluid near the bound is one whose particles scatter so strongly that their mean free path is about their own quantum wavelength: there are no well-defined particles at all, and disturbances are damped as fast as quantum mechanics allows. The quark--gluon plasma flows with almost no friction: the most perfect fluid known.

\sol{strings}{11} If the black hole forms from matter in a pure state, then black hole plus radiation is always pure (if evolution is unitary), and the entanglement entropy of the radiation equals that of the black hole. Early on, each Hawking quantum is entangled with a partner behind the horizon, and the radiation's entropy rises. But the black hole's entanglement entropy cannot exceed its Bekenstein--Hawking entropy, the logarithm of the number of its states, and this shrinks as it evaporates. So the radiation's entropy must turn over when the two curves cross, when the black hole has lost about half its entropy (since $S\propto M^2$, at $M = M_0/\sqrt2$, about $65\%$ of the way through its life). When the black hole is gone the radiation, now everything, is in a pure state, with entropy zero. In Hawking's calculation the radiation is exactly thermal and its entropy rises monotonically to the end, leaving a mixed state from a pure one: a violation of unitarity. In 2019, ``island'' calculations in semiclassical gravity reproduced Page's curve, though how the information gets out remains debated.
